\section{Motivation}

Solving systems of linear equations is the oldest problem of linear algebra — Chinese mathematicians were
row reducing two thousand years ago — and, measured in CPU-hours, it may well be the computation the world
performs most. When engineers simulate a bridge, a chip or tomorrow's weather, the physics is chopped into
millions of small interacting pieces, and the interactions are linear equations; the supercomputer's job
is to solve the resulting system. Concrete examples surround you:

\begin{itemize}
\item \textbf{GPS.} Your position is computed from your distances to several satellites; after a standard
algebraic trick the unknown coordinates satisfy a small linear system, which your phone solves several
times per second.
\item \textbf{Computer graphics and game physics.} Finding where a ray of light meets a surface, or
resolving the forces when a stack of crates topples, means setting up and solving linear systems — many
thousands of them per frame.
\item \textbf{Electronics.} Kirchhoff's laws turn any circuit diagram into a linear system relating its
currents and voltages; circuit simulators such as SPICE are, at their core, linear-system solvers.
\item \textbf{Medical imaging.} A CT scanner measures how much X-ray energy survives along thousands of
lines through the body; reconstructing the image of the inside is solving a huge linear system for the
unknown tissue densities.
\item \textbf{Machine learning.} Fitting a linear model to data by least squares — the workhorse of
statistics — leads to the so-called normal equations: a linear system in the model's parameters.
\end{itemize}

This chapter presents \emph{Gaussian elimination}, the algorithm behind all of the above. It deserves the
same respect as any classic algorithm you know: it is provably correct (elementary row operations preserve
the solution set), it is fast (about $\tfrac23 n^3$ arithmetic operations for $n$ equations with $n$
unknowns), and it is universal. Just as valuable is the \emph{structure theory} that comes with it: a
linear system has either no solution, exactly one, or infinitely many — never, say, exactly two — and the
Kronecker--Capelli theorem decides which case occurs by comparing two ranks. In the infinite case the
solutions are beautifully organised: one particular solution plus a \emph{subspace} of $\K^n$ — the first
place in the course where the abstract theory of Chapter~\ref{ch:vector-spaces} pays off directly.

\section{Systems of linear equations}

\begin{definition}
A \emph{system of $m$ linear equations with $n$ unknowns} $x_1,\ldots,x_n$ is a set of equations
\[
\begin{aligned}
a_{11} x_1 + a_{12} x_2 +\ldots+ a_{1n} x_n &= b_1,\\
a_{21} x_1 + a_{22} x_2 +\ldots+ a_{2n} x_n &= b_2,\\
&\ \vdots \\
a_{m1} x_1 + a_{m2} x_2 +\ldots+ a_{mn} x_n &= b_m,
\end{aligned}
\]
with $a_{ij}, b_i\in\K$. A \emph{solution} of this system is a tuple of values $x_1 = k_1,\ldots, x_n=k_n$
which satisfies all $m$ equalities.
\end{definition}

\begin{example}
The system
\[
\begin{aligned}
3x + y + z &= 1,\\
x + y &= 0
\end{aligned}
\]
has $2$ equations and $3$ unknowns. One checks that $x=0,y=0,z=1$ and $x=-1,y=1,z=3$ are both solutions
(there are others). On the other hand the system
\[
\begin{aligned}
x + y &= 1,\\
x + y &= 0
\end{aligned}
\]
has no solutions, since a solution would imply $0=1$.
\end{example}

\begin{observation}[Matrix representation]
A system of $m$ linear equations with unknowns $x_1,\ldots,x_n$ can be written as a single matrix equation
\[
A\cdot X = B,
\]
where
\[
A = \left( \begin{array}{cccc}
a_{11} & a_{12} & \ldots & a_{1n} \\
a_{21} & a_{22} & \ldots & a_{2n} \\
\vdots & \vdots & \ddots & \vdots \\
a_{m1} & a_{m2} & \ldots & a_{mn}
\end{array} \right), \qquad
X = \left ( \begin{array}{c} x_1 \\ x_2 \\ \vdots \\ x_n \end{array} \right ), \qquad
B = \left ( \begin{array}{c} b_1 \\ b_2 \\ \vdots \\ b_m \end{array} \right ).
\]
\end{observation}

\begin{definition}
Two systems $A_1 X = B_1$ and $A_2 X = B_2$ of $m$ linear equations with $n$ unknowns are \emph{equivalent}
if they have the same solution set.
\end{definition}

\begin{example}
The system $\{\,x - y = 1,\ x + y = 0\,\}$ is equivalent to $\{\,2x - 2y = 2,\ x + y = 0\,\}$.
\end{example}

\begin{definition}
The matrix
\[
(A\mid B) = \left( \begin{array}{cccc|c}
a_{11} & a_{12} & \ldots & a_{1n} & b_1 \\
a_{21} & a_{22} & \ldots & a_{2n} & b_2\\
\vdots & \vdots & \ddots & \vdots & \vdots \\
a_{m1} & a_{m2} & \ldots & a_{mn} & b_m
\end{array} \right)
\]
is called the \emph{augmented matrix} of the system $A\cdot X = B$.
\end{definition}

\begin{example}
The matrix
$\left( \begin{array}{cc|c} 2 & -2 & 2\\ 1 & 1 & 0 \end{array} \right)$
is the augmented matrix of the system $\{\,2x - 2y = 2,\ x + y = 0\,\}$.
\end{example}

\begin{theorem}
Two systems $A_1 X = B_1$ and $A_2 X= B_2$ of $m$ linear equations with $n$ unknowns are equivalent if
their augmented matrices $(A_1\mid B_1)$ and $(A_2\mid B_2)$ are row equivalent.
\end{theorem}

\noindent Recall that two matrices are row equivalent if one can be obtained from the other by elementary row
operations: row switching ($R_i\leftrightarrow R_j$), row scaling by a nonzero scalar
($kR_i\rightarrow R_i$) and row addition ($R_i+k\cdot R_j\rightarrow R_i$).

\section{Solving systems}

\begin{definition}
A system $AX=B$ of linear equations is in \emph{echelon form} if its augmented matrix $(A\mid B)$ is in row
echelon form.
\end{definition}

Since every matrix is row equivalent to a matrix in row echelon form, for any system $AX=B$ there exists an
equivalent system $A'X=B'$ in echelon form, that is, a system in echelon form with the same solution set.
This gives a general method.

\begin{theorem}[Solving a system]
To solve a system $AX=B$:
\begin{itemize}
\item write its augmented matrix $(A\mid B)$;
\item find a matrix $(A'\mid B')$ in row echelon form which is row equivalent to $(A\mid B)$;
\item write the new system $A' X = B'$ and read off its solutions.
\end{itemize}
\end{theorem}

\begin{remark}
This procedure is known as \emph{Gaussian elimination}. Reading the solutions off an echelon system, from
the last equation upwards, is called \emph{back substitution}.
\end{remark}

\begin{definition}
Let $A'X = B'$ be a system in echelon form. The unknowns multiplying the leading entries of the rows of
$A'$ are called \emph{leading variables}; the remaining unknowns are the \emph{free variables}. To read
off the solutions one assigns arbitrary parameter values to the free variables and computes the leading
variables by back substitution.
\end{definition}

\begin{example}[A unique solution]
Consider the system
\[
\begin{aligned}
2x +4y-z &= 11,\\ -4x -3y+3z &= -20,\\ 2x +4y+2z &= 2,
\end{aligned}
\qquad\text{with augmented matrix}\qquad
\left( \begin{array}{ccc|c}
2 & 4 & -1 & 11\\ -4 & -3 & 3 & -20\\ 2 & 4 & 2 & 2
\end{array} \right).
\]
Using elementary row operations we obtain the row echelon form
\[
\left( \begin{array}{ccc|c}
2 & 4 & -1 & 11\\ 0 & 5 & 1 & 2\\ 0 & 0 & 3 & -9
\end{array} \right),
\qquad\text{i.e.}\qquad
\begin{aligned}
2x +4y-z &= 11,\\ 5y+z &= 2,\\ 3z &= -9.
\end{aligned}
\]
The third equation gives $z=-3$; the second then gives $5y-3=2$, so $y=1$; finally $x=2$. The system has the
unique solution $x=2,\ y=1,\ z=-3$.
\end{example}

\begin{example}[Infinitely many solutions]
Consider the system
\[
\begin{aligned}
x + 4y - 3z +2t &= 5,\\ 2x + 8y -5z\phantom{{}+2t} &=12,
\end{aligned}
\qquad\text{with augmented matrix}\qquad
\left( \begin{array}{cccc|c}
1 & 4 & -3 & 2 & 5\\ 2 & 8 & -5 & 0 & 12
\end{array} \right).
\]
By row reduction we obtain
\[
\left( \begin{array}{cccc|c}
1 & 4 & -3 & 2 & 5\\ 0 & 0 & 1 & -4 & 2
\end{array} \right),
\qquad\text{i.e.}\qquad
\begin{aligned}
x + 4y - 3z +2t &= 5,\\ z-4t &= 2.
\end{aligned}
\]
This echelon system has $2$ equations and $4$ unknowns, hence more than one solution. We assign arbitrary
values to the free variables $y=a$ and $t=b$. Substitution gives $z=2+4b$ and $x=11-4a+10b$. The
solution depends on two parameters:
\[
x=11-4a+10b, \quad y = a, \quad z = 2+4b, \quad t=b.
\]
\end{example}

\begin{example}[No solution]
Consider the system
\[
\begin{aligned}
2x +4y &= 0,\\ 2x + 4y &= 1,
\end{aligned}
\qquad\text{with augmented matrix}\qquad
\left( \begin{array}{cc|c} 2 & 4 & 0\\ 2 & 4 & 1 \end{array} \right)
\xrightarrow{\ \ }
\left( \begin{array}{cc|c} 2 & 4 & 0\\ 0 & 0 & 1 \end{array} \right).
\]
The last row reads $0=1$, so this system has no solutions.
\end{example}

\begin{theorem}[Kronecker--Capelli]
A system $AX=B$ of linear equations has a solution if and only if $\rank(A) = \rank(A\mid B)$.
\end{theorem}

\begin{corollary}[Number of solutions]
Let $AX = B$ be a solvable system with $n$ unknowns and let $r = \rank(A) = \rank(A\mid B)$. If $r = n$,
the system has exactly one solution. If $r<n$, the solutions are described by $n-r$ free parameters: over
an infinite field such as $\R$ or $\C$ this means infinitely many solutions, and over a finite field
$\Z_p$ it means exactly $p^{\,n-r}$ of them. In particular, a system over $\R$ has no solution, exactly
one solution, or infinitely many — never, say, exactly two.
\end{corollary}

\begin{example}
For the system from the previous example,
\[
(A\mid B) = \left( \begin{array}{cc|c} 2 & 4 & 0\\ 2 & 4 & 1 \end{array} \right)
\longrightarrow
\left( \begin{array}{cc|c} 2 & 4 & 0\\ 0 & 0 & 1 \end{array} \right),
\]
we read off $\rank(A) = 1$ and $\rank(A\mid B) = 2$. The two ranks differ, which confirms that the system
has no solutions.
\end{example}

\begin{example}[A system with a parameter]
For which values of the parameter $p\in\R$ does the system
\[
\begin{aligned}
x+y+z &= 1,\\ x+2y+3z &= 2,\\ 2x+3y+4z &= p
\end{aligned}
\]
have solutions? Row reducing the augmented matrix,
\[
\left( \begin{array}{ccc|c}
1 & 1 & 1 & 1\\ 1 & 2 & 3 & 2\\ 2 & 3 & 4 & p
\end{array} \right)
\xrightarrow{\substack{R_2-R_1\to R_2\\ R_3-2R_1\to R_3}}
\left( \begin{array}{ccc|c}
1 & 1 & 1 & 1\\ 0 & 1 & 2 & 1\\ 0 & 1 & 2 & p-2
\end{array} \right)
\xrightarrow{R_3-R_2\to R_3}
\left( \begin{array}{ccc|c}
1 & 1 & 1 & 1\\ 0 & 1 & 2 & 1\\ 0 & 0 & 0 & p-3
\end{array} \right).
\]
If $p\neq 3$, the last row reads $0 = p-3\neq 0$ and there are no solutions (equivalently,
$\rank(A) = 2 < 3 = \rank(A\mid B)$). If $p = 3$, then $\rank(A) = \rank(A\mid B) = 2$ and the system has
infinitely many solutions: taking $z = a$ we get $y = 1-2a$ and $x = a$.
\end{example}

\section{Homogeneous systems}

\begin{definition}
A system of linear equations is called \emph{homogeneous} if it has the form $A\cdot X = \0$.
\end{definition}

\noindent A homogeneous system \emph{always} has the solution $x_1=0, x_2=0,\ldots,x_n=0$. Any other solution
(if one exists) is called a \emph{non-trivial} solution.

\begin{fact}
A homogeneous system with fewer equations than unknowns ($m<n$) always has a non-trivial solution: indeed
$\rank(A)\leq m<n$, so after row reduction at least one variable is free.
\end{fact}

\begin{example}
The homogeneous system
\[
\begin{aligned}
x+2y-3z+w &= 0,\\ x-3y +z -2w &=0,\\ 2x+y-3z+5w &= 0
\end{aligned}
\]
has $3$ equations and $4$ unknowns; therefore it has non-trivial solutions.
\end{example}

\begin{example}
The homogeneous system
\[
\begin{aligned}
x+y-z &= 0,\\ 2x-3y+z &=0,\\ x-4y+2z &= 0
\end{aligned}
\qquad\text{reduces to}\qquad
\begin{aligned}
x+y-z &= 0,\\ -5y+3z &= 0.
\end{aligned}
\]
It has a non-trivial solution: assigning $z=a$ gives $y=\tfrac{3}{5}a$ and $x=\tfrac{2}{5}a$.
\end{example}

\begin{example}
The homogeneous system
\[
\begin{aligned}
x+y-z &= 0,\\ 2x+4y-z &=0,\\ 3x+2y+2z &= 0
\end{aligned}
\qquad\text{reduces to}\qquad
\begin{aligned}
x+y-z &= 0,\\ 2y+z &= 0,\\ 11z &= 0.
\end{aligned}
\]
Here the only solution is the trivial one $x=y=z=0$.
\end{example}

\begin{theorem}[Solution space]
Let $AX=\0$ be a system of $m$ homogeneous linear equations with $n$ unknowns. Then the set
$W=\{\vv\in \K^n\mid A\vv=\0\}$ of solutions is a subspace of the vector space $\K^n$. Moreover,
\[
\dim(W) = n-\rank(A).
\]
\end{theorem}

\begin{proof}
We prove that $W$ is a subspace. Take $\vv_1,\vv_2\in W$, so $A\vv_1=\0$ and $A\vv_2=\0$. Then
\[
A(\vv_1+\vv_2) = A\vv_1+A\vv_2 = \0+\0=\0,
\]
hence $\vv_1+\vv_2\in W$. Similarly, for any $k\in \K$ we have $A(k\vv_1)=kA\vv_1=\0$, so $k\cdot \vv_1\in W$.
\end{proof}

\begin{remark}
The dimension formula is visible in the algorithm itself: solving $AX=\0$ by Gaussian elimination leaves
$\rank(A)$ leading variables and $n-\rank(A)$ free variables, and setting one free variable to $1$ and the
others to $0$ — as in the next example — produces a basis of $W$ with $n-\rank(A)$ elements.
\end{remark}

\begin{example}
Consider the homogeneous system
\[
\begin{aligned}
x+2y+3z &= 0,\\ -2x-4y-6z &= 0,
\end{aligned}
\qquad\text{which reduces to}\qquad
x+2y+3z = 0.
\]
The non-leading variables are $y$ and $z$. Putting $y=a$ and $z=b$, the solutions are
\[
\left( \begin{array}{c} x \\ y \\ z \end{array} \right )= \left( \begin{array}{c} -2a-3b \\ a \\ b \end{array} \right ),
\]
so the solution space is
\[
W = \Big\{ \left( \begin{array}{c} -2a-3b \\ a \\ b \end{array} \right ) \;\Big|\; a,b\in\R\Big\}.
\]
Taking $(a,b)=(1,0)$ and $(a,b)=(0,1)$ yields two vectors
\[
\left( \begin{array}{c} -2 \\ 1 \\ 0 \end{array} \right ),\qquad
\left( \begin{array}{c} -3 \\ 0 \\ 1 \end{array} \right )
\]
which form a basis of $W$.
\end{example}

\section{Homogeneous and non-homogeneous systems}

\begin{theorem}
Let $AX=B$ be an arbitrary system of linear equations, let $U$ be its solution set and let $\vv_0\in U$ be a
particular solution. Then
\[
U = \vv_0 + W =\{\vv_0 + \ww \mid \ww\in W\},
\]
where $W$ is the solution space of the associated homogeneous system $AX=\0$.
\end{theorem}

\begin{proof}
Every element of $\vv_0 + W$ is a solution of $AX=B$, since
\[
A(\vv_0 + \ww) = A\vv_0 + A\ww=B+\0=B.
\]
Conversely, if $A\vv=B$, put $\ww=\vv-\vv_0$. Then $\vv = \vv_0+\ww$ and
\[
A\ww = A(\vv-\vv_0) = A\vv-A\vv_0 = B-B = \0,
\]
so $\ww\in W$ and $\vv\in \vv_0+W$.
\end{proof}

\begin{example}
Consider the equation $x+2y+3z = 6$ (a system of one equation with three unknowns). A particular solution
is $\vv_0 = (1,1,1)$, and the associated homogeneous equation $x+2y+3z = 0$ has the solution space
\[
W = \Span{\{(-2,1,0),\ (-3,0,1)\}},
\]
as computed in an earlier example. Hence every solution of $x+2y+3z=6$ has the form
\[
(1,1,1) + a\,(-2,1,0)+b\,(-3,0,1), \qquad a,b\in\R.
\]
\end{example}
